% Generated draft: requires mathematical and source review.
\documentclass[11pt]{article}
\usepackage[margin=25mm]{geometry}
\usepackage{fontspec}
\setmainfont{Latin Modern Roman}
\usepackage{amsmath,amssymb,mathtools,unicode-math}
\setmathfont{Latin Modern Math}
\usepackage{enumitem,needspace,xurl,hyperref}
\setlength{\parindent}{0pt}
\setlength{\parskip}{5pt}
\setlength{\emergencystretch}{3em}
\newcommand{\N}{\mathbb{N}}
\newcommand{\Z}{\mathbb{Z}}
\newcommand{\Q}{\mathbb{Q}}
\newcommand{\R}{\mathbb{R}}
\newcommand{\C}{\mathbb{C}}
\newcommand{\F}{\mathbb{F}}
\newcommand{\A}{\mathbb{A}}
\newcommand{\PP}{\mathbb{P}}
\newcommand{\E}{\mathbb{E}}
\newcommand{\eps}{\varepsilon}
\newcommand{\dd}{\,\mathrm{d}}
\begin{document}
\section*{UnsolvedMath ID: 30006322}
\textbf{Source problem code:} catalogue entry
\subsection{Definitions and mathematical statement}
Write \(G^\vee\) for the dual group appearing in the intended local geometric duality, \(\mathfrak g^\vee=\operatorname{Lie}(G^\vee)\), and \(\mathfrak g^{\vee *}=(\mathfrak g^\vee)^*\), equipped with the coadjoint action. The source does not specify the base field, the hypotheses on the group, or the precise version of local geometric duality.

For any affine Hamiltonian \(G^\vee\)-space \(M^\vee\) in that setting, let
\[
\mu:M^\vee\longrightarrow\mathfrak g^{\vee *}
\]
be its equivariant moment map. Smoothness is not imposed by the supplied excerpt, which explicitly discusses potentially nonsmooth Hamiltonian spaces. Write \(\operatorname{QCoh}(X)^{G^\vee}\) for the category of equivariant quasi-coherent sheaves on a \(G^\vee\)-space \(X\), with any derived or grading conventions to be specified by the intended duality. The excerpt describes the adjunction
\[
\mu^*:\operatorname{QCoh}(\mathfrak g^{\vee *})^{G^\vee}
\rightleftarrows
\operatorname{QCoh}(M^\vee)^{G^\vee}:\mu_*.
\]
The corresponding function-algebra object is
\[
\mathcal A_{M^\vee}:=\mu_*\mathcal O_{M^\vee}
\quad\text{in}\quad
\operatorname{QCoh}(\mathfrak g^{\vee *})^{G^\vee}.
\]
Here \(\mathcal O_{M^\vee}\) denotes the structure sheaf. In the ordinary affine interpretation, its pushforward represents the full algebra \(\mathcal O(M^\vee)=\Gamma(M^\vee,\mathcal O_{M^\vee})\), retaining its \(G^\vee\)-action and its algebra structure over \(\mathcal O(\mathfrak g^{\vee *})\) induced by \(\mu\). This distinguishes the full equivariant algebra from the invariant subalgebra \(\mathcal O(M^\vee)^{G^\vee}\).

\textbf{Requested conjectural upgrade.} Does the local geometric duality in the cited setting admit an equivariant refinement of its Borel--Moore homology algebra description that recovers \(\mathcal A_{M^\vee}\), and hence the full function algebra with its equivariant structure, before taking \(G^\vee\)-invariants? The supplied curation calls this the Plancherel-algebra conjecture; the excerpt points to reference [3, Conjecture 8.1.8].

\textbf{Unresolved scope.} The supplied excerpt omits the preceding Borel--Moore isomorphism and ends before specifying the adjunction corresponding to \(\mu^*\dashv\mu_*\). It invokes compatibility with pointings and linearity over a category denoted \(H_G\), identified there with an equivariant quasi-coherent category, but does not provide the necessary definitions or grading conventions. Consequently the geometric-side algebra, the precise comparison, and the class of spaces for which the upgrade is conjectured cannot be stated completely from this record. The question above records the intended target without asserting an unspecified equivalence.
\subsection{Short English statement}
The question asks whether the cited local geometric duality can recover the entire function algebra of the dual Hamiltonian space, including its group action and moment-map structure, rather than only its invariant functions. The source omits the geometric-side construction and several hypotheses, so this is a qualified statement-only draft requiring mathematical review.
\subsection{Sources}
\begin{itemize}
\item Supplied local UnsolvedMath registry record, permanent ID \(30006322\), titled ``Removing Dual-Group Invariants from Plancherel Geometry''; exact source problem code \texttt{catalogue entry}. The record attributes the problem to the 2025 Oberwolfach report and supplies no individual proposer.
\item \emph{Arbeitsgemeinschaft: Relative Langlands Duality} (2025), Oberwolfach Reports, pp. 811--882, DOI \(10.4171/owr/2025/18\). Supplied source and literature link: \url{https://doi.org/10.4171/owr/2025/18}.
\item \emph{Relative Langlands Duality}. Supplied literature link: \url{https://arxiv.org/abs/2409.04677}. The excerpt cites [3, Conjecture 8.1.8], but the supplied record does not reproduce the bibliography entry identifying [3].
\item Literature assessment dated exactly \texttt{2026-08-22}: the supplied curation labels the problem open, identifies the proposed upgrade as the Plancherel-algebra conjecture, and reports that the primary paper and the 2025 report formulate it conjecturally, with no general resolution located. This is supplied dated curation, not an independently verified literature assessment; the record requests independent verification of time-sensitive status and cited claims.
\item The registry labels its statement \texttt{corrected\_verified} and reports that the official report was checked, the question isolated from surrounding exposition, and notation repaired. These are supplied source-verification claims, not verification performed for this draft.
\end{itemize}
\end{document}
